\lecture{35}{Алгоритмы распределённой оптимизации}

\sourcevideo
    {https://www.youtube.com/playlist?list=PLOzRYVm0a65fhKDS8cYIC7YCuVGJ4FOJx}
    {Week 10: Lecture 35}

Пусть \(N\) агентов связаны коммуникационным графом, а агент \(i\)
знает только собственную функцию
\(f_i\colon\mathbb R^d\to\mathbb R\). Требуется минимизировать сумму
локальных функций, не собирая их в одном узле. Основной метод лекции
--- распределённый градиентный спуск (distributed gradient descent,
DGD), сочетающий локальное усреднение и градиентный шаг.

\section{Постановка задачи}

Централизованная задача имеет вид
\[
    \min_{x\in\mathbb R^d} F(x),
    \qquad
    F(x)=\sum_{i=1}^{N}f_i(x).
\]
При доступе ко всем функциям обычный градиентный спуск записывается как
\[
    x^{k+1}
    =
    x^k-\eta_k\nabla F(x^k),
    \qquad
    \nabla F(x)=\sum_{i=1}^{N}\nabla f_i(x).
\]
В распределённой сети агент не может вычислить эту сумму. Он хранит
локальную копию \(x_i\in\mathbb R^d\), вычисляет только
\(\nabla f_i(x_i)\) и обменивается сообщениями с соседями.
Эквивалентная формулировка имеет вид
\[
\begin{aligned}
    \min_{x_1,\ldots,x_N}\quad
    &\sum_{i=1}^{N}f_i(x_i),
    \\
    \text{при условии}\quad
    &x_1=\cdots=x_N.
\end{aligned}
\]
Ограничение консенсуса возвращает одну общую переменную и тем самым
восстанавливает исходную задачу.

Для примера рассмотрим
\[
    f_i(x)=\frac12(x-i)^2,
    \qquad i=1,\ldots,5.
\]
Тогда
\[
    F'(x)=\sum_{i=1}^{5}(x-i)=5x-15,
\]
поэтому глобальный минимум равен \(x^\star=3\). Изолированный агент
получил бы локальный минимум \(x_i^\star=i\), так что без коммуникации
глобальное решение недостижимо.

\section{Матрица смешивания}

Обмен сообщениями задаётся связным неориентированным графом
\[
    \mathcal G=(\mathcal V,\mathcal E),
    \qquad
    \mathcal V=\{1,\ldots,N\},
\]
и матрицей смешивания
\(W=(w_{ij})\in\mathbb R^{N\times N}\). Для базового анализа
предполагается
\[
    w_{ij}\ge0,
    \qquad
    w_{ij}=0
    \quad\text{при }j\notin\mathcal N_i\cup\{i\},
\]
а также двойная стохастичность
\[
    W\mathbf1=\mathbf1,
    \qquad
    \mathbf1^{\mathsf T}W=\mathbf1^{\mathsf T}.
\]
На неориентированном графе этим требованиям удовлетворяют, например,
симметричные веса Метрополиса:
\[
    w_{ij}
    =
    \frac{1}{1+\max\{d_i,d_j\}}
\]
для соседних \(i\ne j\) и
\[
    w_{ii}
    =
    1-\sum_{j\in\mathcal N_i}w_{ij}.
\]
Строчная стохастичность делает локальное обновление выпуклой
комбинацией соседних состояний, а столбцовая сохраняет их среднее.

\section{Распределённый градиентный спуск}

На итерации \(k\) агент сначала смешивает локальные переменные,
\[
    x_i^{k+\frac12}
    =
    \sum_{j=1}^{N}w_{ij}x_j^k,
\]
а затем выполняет градиентный шаг по своей функции:
\[
    x_i^{k+1}
    =
    x_i^{k+\frac12}
    -
    \eta_k\nabla f_i\!\left(x_i^{k+\frac12}\right).
\]
Смешивание уменьшает рассогласование, тогда как локальные градиенты
двигают агентов к минимумам разных функций. Возможен и обратный
порядок:
\[
    x_i^{k+\frac12}
    =
    x_i^k-\eta_k\nabla f_i(x_i^k),
    \qquad
    x_i^{k+1}
    =
    \sum_{j=1}^{N}w_{ij}x_j^{k+\frac12}.
\]
Обе схемы сочетают оптимизационный и консенсусный шаги, но имеют
разные рекуррентные соотношения и требуют отдельных доказательств.

Обозначим среднее состояние через
\[
    \bar x^k
    =
    \frac1N\sum_{i=1}^{N}x_i^k.
\]
Для первого порядка шагов столбцовая стохастичность даёт
\[
\begin{aligned}
    \frac1N\sum_{i=1}^{N}x_i^{k+\frac12}
    &=
    \frac1N\sum_{j=1}^{N}
    \left(\sum_{i=1}^{N}w_{ij}\right)x_j^k
    \\
    &=
    \bar x^k.
\end{aligned}
\]
После локального градиентного шага
\[
    \bar x^{k+1}
    =
    \bar x^k
    -
    \frac{\eta_k}{N}
    \sum_{i=1}^{N}
    \nabla f_i\!\left(x_i^{k+\frac12}\right).
\]
Если
\[
    \bar F(x)=\frac1N\sum_{i=1}^{N}f_i(x),
\]
то централизованный шаг использовал бы
\(\nabla\bar F(\bar x^k)\). Различие возникает потому, что локальные
градиенты вычисляются в разных точках.

\section{Средняя динамика и рассогласование}

Добавляя и вычитая градиенты в средней точке, получаем
\[
    \bar x^{k+1}
    =
    \bar x^k
    -
    \eta_k\nabla\bar F(\bar x^k)
    +
    e^k,
\]
где
\[
    e^k
    =
    \frac{\eta_k}{N}
    \sum_{i=1}^{N}
    \left[
        \nabla f_i(\bar x^k)
        -
        \nabla f_i\!\left(x_i^{k+\frac12}\right)
    \right].
\]
Следовательно, среднее состояние подчиняется возмущённому
централизованному градиентному спуску. Если каждый градиент
\(L_i\)-липшицев и \(L=\max_iL_i\), то
\[
    \lVert e^k\rVert
    \le
    \frac{\eta_kL}{N}
    \sum_{i=1}^{N}
    \left\lVert x_i^{k+\frac12}-\bar x^k\right\rVert.
\]
Ошибка исчезает, когда шаг и рассогласование становятся достаточно
малыми.

Для матричного описания положим
\[
    \mathbf x^k
    =
    \begin{pmatrix}
        x_1^k\\
        \vdots\\
        x_N^k
    \end{pmatrix},
    \qquad
    J=\frac1N\mathbf1\mathbf1^{\mathsf T}.
\]
Консенсусная часть равна
\(\mathbf1\otimes\bar x^k\), а рассогласование ---
\[
    \mathbf x^k-\mathbf1\otimes\bar x^k.
\]
Для симметричной примитивной дважды стохастической матрицы
\[
    \rho
    =
    \lVert W-J\rVert_2
    <1.
\]
Чистое смешивание сжимает рассогласование с коэффициентом не больше
\(\rho\), однако следующий локальный градиентный шаг снова его создаёт.
Поэтому анализ DGD должен одновременно контролировать оптимизационную
ошибку среднего состояния и расстояние локальных переменных до
консенсусного подпространства.

\section{Сходимость и выбор шага}

Классический анализ точной сходимости использует убывающий шаг
\[
    \eta_k>0,
    \qquad
    \sum_{k=0}^{\infty}\eta_k=\infty,
    \qquad
    \sum_{k=0}^{\infty}\eta_k^2<\infty.
\]
Первое условие не позволяет суммарному движению оказаться конечным до
достижения решения, а второе делает квадратичные возмущения
суммируемыми. Типичный выбор --- \(\eta_k=\eta_0/(k+1)\).

При связном неориентированном графе, примитивной дважды
стохастической матрице \(W\), выпуклых гладких функциях, непустом
множестве решений и стандартных дополнительных условиях
ограниченности можно доказать
\[
    \lVert x_i^k-x_j^k\rVert\to0
\]
и
\[
    \operatorname{dist}
    \left(
        \bar x^k,
        \operatorname*{arg\,min}F
    \right)
    \to0.
\]
Если минимизатор единственен, то
\[
    x_i^k\to x^\star
\]
для всех агентов.

При постоянном шаге \(\eta_k=\eta>0\) обычный DGD в общем случае
сходится лишь к окрестности решения. Её размер уменьшается вместе с
\(\eta\) и зависит от спектральных свойств графа и различия локальных
градиентов. Для точной сходимости с постоянным шагом вводят
корректирующие механизмы. Упомянутый в лекции алгоритм EXTRA устраняет
стационарное смещение DGD и при подходящих предположениях сходится к
точному решению; его формула и доказательство здесь не выводятся.

\section{Коммуникационная стоимость и ограничения}

На каждой итерации агент передаёт соседям вектор
\(x_i^k\in\mathbb R^d\). Поэтому суммарный объём обмена по активным
рёбрам имеет порядок
\[
    O\bigl(|\mathcal E|d\bigr),
\]
а локальная вычислительная стоимость определяется вычислением
\(\nabla f_i(x_i)\). Передача дополнительных градиентных или
инерционных переменных сохраняет порядок \(d\) на сообщение, но
увеличивает постоянный множитель.

Базовая схема предполагает синхронные итерации, фиксированный связный
граф, точную передачу сообщений и отсутствие злонамеренных агентов.
Для ориентированных графов двойная стохастичность может быть
недоступна, для изменяющейся топологии требуется совместная связность,
а задержки, асинхронность и стохастические градиенты приводят к другим
оценкам. Следующая лекция вводит градиентный трекинг, позволяющий
устранить смещение постоянного шага за счёт дополнительной локальной
переменной.
